\documentclass[11pt, a4paper]{article}

\usepackage{amsmath, amssymb, amsthm}
\usepackage{algorithm}
\usepackage{algpseudocode}
\usepackage{booktabs}
\usepackage{graphicx}
\usepackage{hyperref}
\usepackage{geometry}
\usepackage{enumitem}
\usepackage{xcolor}
\usepackage{tcolorbox}

% Readable blue for highlighted sections (subsections 4.2, 5.1; sections 6, 8--11)
\newcommand{\bluecontentbegin}{\begingroup\color{blue!70!black}}
\newcommand{\bluecontentend}{\endgroup}

\geometry{margin=1in}

\newtheorem{definition}{Definition}
\newtheorem{remark}{Remark}

\title{Runtime Prediction for Heterogeneous Serverless Scheduling:\\
    From Benchmark-Derived Scaling Factors to Adaptive Prediction}
\author{Serverless Scheduler Project}
\date{\today}

\begin{document}

\maketitle

\begin{abstract}
We present a runtime prediction model for scheduling fixed-workload serverless
functions across a cluster of heterogeneous machines. The model requires no
historical data for a given (function, machine) pair: it derives predictions
from offline benchmarks of functions and machines independently. As historical
observations accumulate, the model smoothly transitions to an exponential
moving average of observed runtimes. We describe the mathematical framework,
the benchmarking methodology, and the complete prediction algorithm.
\end{abstract}

\tableofcontents
\newpage

% =============================================================================
\section{Problem Setting}
% =============================================================================

We have a fixed set of serverless functions $\mathcal{F} = \{f_1, \ldots, f_N\}$
and a heterogeneous cluster of machines $\mathcal{M} = \{m_1, \ldots, m_K\}$.
A scheduler receives batches of function invocation requests and must assign
each invocation to a machine so as to minimize total batch completion time.
The quality of this assignment depends critically on the accuracy of the
predicted runtime $\hat{t}(f, m)$ for every (function, machine) pair.

\subsection{Key Assumptions}

\begin{enumerate}[label=(\roman*)]
    \item \textbf{Fixed workload.} Each function performs the same computation
          on every invocation. Input size does not vary. Therefore, the runtime
          of function $f$ on machine $m$ is essentially deterministic, subject
          only to system noise (OS scheduling jitter, cache state, background
          load).
    \item \textbf{Sparse history.} Most $(f, m)$ combinations will not have
          been observed before. With $N$ functions and $K$ machines, the space
          of pairs is $N \times K$, but only a small fraction will have
          historical data at any point.
    \item \textbf{Benchmarking is affordable.} We can afford to run each
          function multiple times on a dedicated Benchmarking Environment
          Machine (BEM), and to run standardized hardware benchmarks on each
          cluster machine.
\end{enumerate}

\subsection{Notation}

\begin{center}
\begin{tabular}{cl}
    \toprule
    Symbol & Meaning \\
    \midrule
    $f \in \mathcal{F}$ & A serverless function \\
    $m \in \mathcal{M}$ & A cluster machine (provider) \\
    $m_0$ & The Benchmarking Environment Machine (BEM) \\
    $t_{\mathrm{ref}}(f)$ & Reference runtime of $f$ on BEM \\
    $\hat{t}(f, m)$ & Predicted runtime of $f$ on $m$ \\
    $S_i(m)$ & Score of $m$ on resource dimension $i$ \\
    $w_i(f)$ & Sensitivity weight of $f$ to dimension $i$ \\
    $r_i(m)$ & Performance ratio $S_i(m_0) / S_i(m)$ \\
    $n(f, m)$ & Number of historical observations for the pair \\
    $\mathrm{EMA}(f, m)$ & Exponential moving average of observed runtimes \\
    \bottomrule
\end{tabular}
\end{center}

% =============================================================================
\section{Current Implementation and Its Limitations}
\label{sec:current}
% =============================================================================

Before presenting the proposed model, we document how runtime prediction
currently works end-to-end, and where it breaks.

\subsection{End-to-End Prediction Flow (Status Quo)}

When the scheduler receives a batch of function invocations, it builds a cost
matrix $C[m][f]$ of predicted runtimes for every (provider, service) pair.
This proceeds in three passes:

\paragraph{Pass 1: Historical lookup
(\texttt{get\_predicted\_runtimes\_db\_pass}).}
For each $(f, m)$, query the \texttt{Job} table for the single most recent
completed job with that \texttt{(provider\_id, service\_id)} pair and return
its \texttt{run\_time}. If no such row exists, mark the service as
``needing prediction.''

\paragraph{Pass 2: MQTT prediction request
(\texttt{\_fetch\_predicted\_runtimes\_mqtt\_batch}).}
For every service that had no historical data, the scheduler publishes a
\texttt{PREDICT\_REQUEST} message over MQTT to the relevant provider. The
provider receives this, calls \texttt{trainAndPredict}, and replies with
\texttt{PREDICT\_RESPONSE} containing predicted milliseconds. The scheduler
blocks (up to a 5\,s timeout) waiting for all responses.

On the provider side, \texttt{trainAndPredict} works as follows:
\begin{enumerate}[nosep]
    \item Load \emph{all} past training rows from
          \texttt{TrainingData/eff\_score\_data.txt} --- these are not
          function-specific; every function's runs are mixed together.
    \item Build features $X = [\texttt{cpu\_usage} \times
          \texttt{cpu\_eff},\; \texttt{mem\_usage} \times
          \texttt{mem\_eff}]$ and target $y = \texttt{actual\_runtime}$.
    \item Fit a fresh \texttt{sklearn.linear\_model.LinearRegression}
          from scratch.
    \item For the requested service, load its reference CPU/memory from
          \texttt{Reference\_Provider\_Data.txt}, multiply by the
          provider's global efficiency scores, and call
          \texttt{model.predict}.
\end{enumerate}

\paragraph{Pass 3: Pull time adjustment
(\texttt{get\_predicted\_runtimes\_cache\_pass}).}
Adjust every predicted runtime based on container image cache state: if
not cached, add the full last observed \texttt{pull\_time}; if cached on
disk, add $0.2 \times \texttt{pull\_time}$; if in memory, add nothing.
If anything fails at any point, fall back to a hard-coded default of
1000\,ms.

The specific deficiencies of this pipeline are enumerated in
Appendix~\ref{app:broken}.

% =============================================================================
\section{Resource Dimensions}
\label{sec:resources}
% =============================================================================

We model runtime as being influenced by four resource dimensions:

\begin{definition}[Resource Dimensions]
\[
    \mathcal{R} = \{\, \texttt{cpu},\; \texttt{mem},\; \texttt{disk},\; \texttt{net} \,\}
\]
\begin{enumerate}[label=(\alph*)]
    \item \textbf{CPU} --- processor compute throughput (cycles/second, or a
          composite score from a standardized benchmark such as \texttt{sysbench}).
    \item \textbf{Memory} --- memory bandwidth (MB/s, as measured by a STREAM
          or equivalent benchmark).
    \item \textbf{Disk I/O} --- sequential read/write throughput (MB/s, as
          measured by \texttt{fio} or equivalent).
    \item \textbf{Network I/O} --- network bandwidth (Mbps, as measured by
          \texttt{iperf3} to a fixed reference endpoint).
\end{enumerate}
\end{definition}

\begin{remark}
The current codebase tracks only CPU and memory efficiency scores
(\texttt{cpu\_efficiency\_score}, \texttt{memory\_efficiency\_score} in the
\texttt{User} model). Disk and network dimensions are proposed additions.
\end{remark}

% =============================================================================
\section{Benchmarking Methodology}
\label{sec:benchmarking}
% =============================================================================

% -----------------------------------------------------------------------------
\subsection{Machine Benchmarking}
\label{sec:machine-bench}
% -----------------------------------------------------------------------------

For each machine $m \in \mathcal{M} \cup \{m_0\}$, we run standardized
benchmarks to obtain absolute performance scores:

\[
    S_{\texttt{cpu}}(m), \quad S_{\texttt{mem}}(m), \quad
    S_{\texttt{disk}}(m), \quad S_{\texttt{net}}(m)
\]

These are hardware-intrinsic and need only be measured once per machine (or
re-measured if hardware changes). From these, we compute the
\emph{performance ratio} of each machine relative to the BEM:

\begin{equation}
\label{eq:ratio}
    r_i(m) = \frac{S_i(m_0)}{S_i(m)}, \qquad i \in \mathcal{R}
\end{equation}

\begin{itemize}
    \item $r_i(m) = 1$ when $m$ matches the BEM on dimension $i$.
    \item $r_i(m) > 1$ when $m$ is \emph{slower} than the BEM (runtime scales up).
    \item $r_i(m) < 1$ when $m$ is \emph{faster} than the BEM (runtime scales down).
\end{itemize}

For the BEM itself, $r_i(m_0) = 1$ for all $i$, serving as the identity anchor.

% -----------------------------------------------------------------------------
\subsection{Function Benchmarking on BEM}
\label{sec:func-bench}
% -----------------------------------------------------------------------------

For each function $f \in \mathcal{F}$, we perform two stages of benchmarking
on the BEM.

\subsubsection{Stage 1: Reference Runtime}

Run function $f$ on $m_0$ a total of $B$ times (we recommend $B \geq 5$)
under normal (unthrottled) conditions. Record the runtimes
$\{t_1, \ldots, t_B\}$ and take the median:

\begin{equation}
\label{eq:tref}
    t_{\mathrm{ref}}(f) = \mathrm{median}(t_1, \ldots, t_B)
\end{equation}

The median is preferred over the mean for robustness to outliers (cold-start
effects, GC pauses, etc.).

\subsubsection{Stage 2: Resource Sensitivity via Throttling}
\label{sec:throttling}

To determine how sensitive function $f$ is to each resource dimension, we
throttle one resource at a time on the BEM and measure the resulting slowdown.

For each resource $i \in \mathcal{R}$, throttle resource $i$ to a fraction
$\theta$ of its full capacity (e.g., $\theta = 0.5$) while leaving all other
resources unthrottled. Run $f$ a total of $B'$ times ($B' \geq 3$) and
record the median throttled runtime $t_i^{\mathrm{thr}}(f)$.

\begin{remark}[Throttling Mechanisms]
Practical throttling on the BEM can be achieved via:
\begin{itemize}
    \item \textbf{CPU:} Docker \texttt{--cpus} flag or \texttt{cgroup} CPU quota.
    \item \textbf{Memory bandwidth:} Synthetic memory pressure via
          \texttt{stress-ng --vm}, or \texttt{cgroup} memory bandwidth limits
          where supported.
    \item \textbf{Disk I/O:} Docker \texttt{--device-read-bps} /
          \texttt{--device-write-bps} or \texttt{cgroup} \texttt{blkio} limits.
    \item \textbf{Network:} Linux \texttt{tc} (traffic control) with
          \texttt{netem} or \texttt{tbf} queueing disciplines.
\end{itemize}
\end{remark}

\begin{remark}[Choosing the Throttle Fraction $\theta$]
\label{rem:theta}
The throttle fraction should reflect the performance range of the actual
cluster. If the weakest machine has $0.3\times$ the CPU throughput of the BEM,
throttling to $\theta = 0.3$ during weight discovery produces deviations that
are calibrated to the actual prediction range. A uniform $\theta = 0.5$ across
all resources is a reasonable default.
\end{remark}

Compute the \emph{relative deviation} for each resource:

\begin{equation}
\label{eq:deviation}
    d_i(f) = \frac{t_i^{\mathrm{thr}}(f) - t_{\mathrm{ref}}(f)}{t_{\mathrm{ref}}(f)},
    \qquad i \in \mathcal{R}
\end{equation}

This measures what fraction of additional runtime is attributable to degrading
resource $i$. A CPU-bound function will exhibit large $d_{\texttt{cpu}}$ and
small $d_{\texttt{mem}}$, $d_{\texttt{disk}}$, $d_{\texttt{net}}$.

\paragraph{Clamping.}
In rare cases, throttling a resource may not slow down the function (or may
even speed it up due to second-order effects such as altered garbage collection
behavior). We clamp:

\begin{equation}
\label{eq:clamp}
    d_i(f) \leftarrow \max\bigl(0,\; d_i(f)\bigr)
\end{equation}

\paragraph{Normalization to Weights.}

\begin{equation}
\label{eq:weights}
    w_i(f) = \frac{d_i(f)}{\displaystyle\sum_{j \in \mathcal{R}} d_j(f)},
    \qquad \text{so that} \quad \sum_{i \in \mathcal{R}} w_i(f) = 1
\end{equation}

\paragraph{Fallback for fast functions.}
If $\sum_j d_j(f) < \epsilon$ (the function is too fast for throttling to
produce measurable deviations), assign equal weights:
$w_i(f) = 1/|\mathcal{R}| = 0.25$ for all $i$.

% =============================================================================
\section{The Scaling Factor Model}
\label{sec:scaling}
% =============================================================================

Given the benchmark-derived quantities from Sections
\ref{sec:machine-bench}--\ref{sec:func-bench}, we can predict the runtime of
any $(f, m)$ pair without any historical observation of that specific pair.

\begin{definition}[Scaling Factor]
\label{def:scaling}
The scaling factor for function $f$ on machine $m$ relative to the BEM is:

\begin{equation}
\label{eq:scaling}
    \boxed{\;\sigma(f, m) = \sum_{i \in \mathcal{R}} w_i(f) \cdot r_i(m)\;}
\end{equation}

where $w_i(f)$ are the function's resource sensitivity weights
(Equation~\ref{eq:weights}) and $r_i(m)$ are the machine's performance ratios
(Equation~\ref{eq:ratio}).
\end{definition}

The cold-start predicted runtime is then:

\begin{equation}
\label{eq:coldstart}
    \hat{t}_{\mathrm{cold}}(f, m) = t_{\mathrm{ref}}(f) \cdot \sigma(f, m)
\end{equation}

\subsection{Intuition}

The scaling factor is a weighted average of per-resource slowdown/speedup
ratios relative to the BEM, where the weights encode how much the function
cares about each resource. Consider two extremes:

\begin{itemize}
    \item A purely CPU-bound function ($w_{\texttt{cpu}} = 1$, others $= 0$):
          $\sigma = r_{\texttt{cpu}}(m)$. The prediction scales linearly with
          the CPU performance ratio. If the target has $2\times$ the CPU
          throughput, the predicted runtime halves.
    \item A function equally sensitive to all resources ($w_i = 0.25$):
          $\sigma$ is the unweighted mean of all four ratios. The prediction
          reflects the average hardware capability.
\end{itemize}

\subsection{Additive vs.\ Multiplicative Formulation}
\bluecontentbegin

An alternative multiplicative formulation is:

\begin{equation}
\label{eq:multiplicative}
    \sigma_{\mathrm{mult}}(f, m)
    = \prod_{i \in \mathcal{R}} r_i(m)^{\,w_i(f)}
\end{equation}

The \textbf{additive} model (Equation~\ref{eq:scaling}) is appropriate when the
function executes resource-dependent phases \emph{sequentially} --- e.g., a
compute phase followed by a network transfer phase. The total runtime is the
sum of phase durations, each scaling independently.

The \textbf{multiplicative} model (Equation~\ref{eq:multiplicative}) is
appropriate when resources interact concurrently --- e.g., a function that
simultaneously streams data from network while computing.

For serverless functions that are typically short-lived and exhibit sequential
resource access patterns, \textbf{the additive model is the recommended
default}. We note that both formulations agree when all ratios are equal
($r_i = c$ for all $i$), since $\sum w_i \cdot c = c = \prod c^{w_i}$.
\bluecontentend

% =============================================================================
\section{Historical Data: Exponential Moving Average}
\label{sec:ema}
% =============================================================================

When a $(f, m)$ pair has been observed before, the observed runtime is a more
direct signal than the benchmark-derived prediction. We maintain an exponential
moving average (EMA) of observed runtimes, updated upon each job completion.

\begin{definition}[EMA Update]
\label{def:ema}
Let $x_n$ be the observed runtime of the $n$-th execution of function $f$ on
machine $m$. The EMA is defined recursively:

\begin{equation}
\label{eq:ema}
    \mathrm{EMA}_n(f, m) =
    \begin{cases}
        \hat{t}_{\mathrm{cold}}(f, m) & \text{if } n = 0 \text{ (initialization)} \\[4pt]
        \alpha \cdot x_n + (1 - \alpha) \cdot \mathrm{EMA}_{n-1}(f, m) & \text{if } n \geq 1
    \end{cases}
\end{equation}

where $\alpha \in (0, 1)$ is the smoothing factor. We recommend
$\alpha \in [0.2, 0.3]$.
\end{definition}

\paragraph{Choice of $\alpha$.}
A smaller $\alpha$ (e.g., $0.1$) produces smoother estimates that are slow to
react to regime changes (e.g., provider load shifts). A larger $\alpha$ (e.g.,
$0.4$) reacts quickly but is more susceptible to noise. For fixed-workload
functions on dedicated infrastructure, $\alpha = 0.2$ balances stability and
responsiveness.

\paragraph{Storage.}
The EMA requires storing a single floating-point value per observed $(f, m)$
pair, plus a count $n(f, m)$. This is an $O(1)$ update --- no historical
data needs to be fetched or retained.

\paragraph{Initialization.}
The EMA is initialized to the cold-start prediction
$\hat{t}_{\mathrm{cold}}(f, m)$, ensuring continuity when the first
observation arrives.

\subsection{Implementation Note}
\bluecontentbegin

In the current codebase, the \texttt{Job} model stores \texttt{run\_time} per
completed job, and \texttt{get\_latest\_run\_time} returns only the most recent
value. The EMA replaces this single-sample lookup with a maintained aggregate.
This requires either:

\begin{enumerate}[label=(\alph*)]
    \item A new field \texttt{ema\_runtime} on a \texttt{(provider, service)}
          summary table (or on the \texttt{Job} model itself), updated on each
          \texttt{JOB\_COMPLETE} event. \textbf{(Recommended.)}
    \item A new model, e.g., \texttt{PairRuntimeStats}, with fields:
          \texttt{provider\_id}, \texttt{service\_id},
          \texttt{ema\_runtime}, \texttt{observation\_count}.
\end{enumerate}
\bluecontentend

% =============================================================================
\section{Adaptive Blending: Cold-Start to Steady-State}
\label{sec:blending}
% =============================================================================
\bluecontentbegin

The two prediction sources --- the benchmark-derived scaling factor and the
EMA of observations --- are complementary. The scaling factor is always
available but approximate; the EMA is precise but requires observations. We
blend them with a \emph{confidence-weighted} scheme.

\begin{definition}[Blended Prediction]
\label{def:blended}
\begin{equation}
\label{eq:blended}
    \boxed{\;
    \hat{t}(f, m) = \frac{n}{n + \kappa} \cdot \mathrm{EMA}(f, m)
                   + \frac{\kappa}{n + \kappa} \cdot \hat{t}_{\mathrm{cold}}(f, m)
    \;}
\end{equation}

where $n = n(f, m)$ is the number of historical observations and
$\kappa > 0$ is a tuning constant that controls how quickly we trust
observations over benchmarks.
\end{definition}

\paragraph{Behavior.}
\begin{itemize}
    \item $n = 0$: $\hat{t} = \hat{t}_{\mathrm{cold}}$. Pure scaling factor.
    \item $n = \kappa$: $\hat{t}$ is a 50/50 blend.
    \item $n \gg \kappa$: $\hat{t} \to \mathrm{EMA}$. History dominates.
\end{itemize}

We recommend $\kappa \in [3, 7]$. With $\kappa = 5$, after 5 observations the
prediction is already $50\%$ driven by actual data; after 15 observations it is
$75\%$ driven by data.

\begin{remark}
This blending formula is formally equivalent to Bayesian shrinkage with a
Gaussian prior (the scaling factor prediction) and Gaussian likelihood (the
observed runtimes). The parameter $\kappa$ plays the role of the prior
precision relative to the data precision. We achieve the effect of Bayesian
updating without any explicit probabilistic machinery.
\end{remark}
\bluecontentend

% =============================================================================
\section{Image Pull Time Adjustment}
\label{sec:pulltime}
% =============================================================================

Total predicted latency adds a pull-time component:
$\hat{t}_{\mathrm{total}}(f, m) = \hat{t}(f, m) + \hat{t}_{\mathrm{pull}}(f, m)$.

Let $V(f)$ be the compressed image size (MB), recorded at function
registration from the image manifest. $S_{\texttt{net}}(m)$ and
$S_{\texttt{disk}}(m)$ are the same machine benchmark scores from
Section~\ref{sec:machine-bench}, in MB/s.

\begin{equation}
\label{eq:pull-piecewise}
    \hat{t}_{\mathrm{pull}}(f, m) =
    \begin{cases}
        0 & \text{memory cache} \\[6pt]
        V(f)\;/\;S_{\texttt{disk}}(m) & \text{disk cache} \\[6pt]
        V(f)\;/\;S_{\texttt{net}}(m) & \text{cold pull}
    \end{cases}
\end{equation}

\paragraph{Why the current heuristic is baseless.}
The codebase adds \texttt{int(pull\_time * 0.2)} for disk-cached images.
The constant $0.2$ has no derivation: it does not depend on $V(f)$,
does not vary across machines with different disk speeds, and reuses a
single past measurement that conflates registry RTT, congestion, and
load. Equation~\ref{eq:pull-piecewise} replaces this with image size
divided by benchmarked throughput --- the same principle as the compute
scaling factor.

% =============================================================================
% APPENDICES
% =============================================================================
\newpage
\appendix
\begin{center}
    {\Large\bfseries Appendices}
\end{center}
\addcontentsline{toc}{section}{Appendices}
\bigskip

% =============================================================================
\section{Complete Prediction Algorithm}
\label{sec:algorithm}
% =============================================================================
\bluecontentbegin

\begin{algorithm}[H]
\caption{Predict Runtime for $(f, m)$}
\label{alg:predict}
\begin{algorithmic}[1]
\Require Function $f$, Machine $m$
\Require Preloaded: $t_{\mathrm{ref}}(f)$,
         $\mathbf{w}(f) = [w_i(f)]$,
         $\mathbf{r}(m) = [r_i(m)]$
\Require Stored: $\mathrm{EMA}(f,m)$, $n(f,m)$
         (if any observations exist)
\Require Parameters: $\kappa$

\Statex

\Function{PredictRuntime}{$f, m$}
    \State $\sigma \gets \displaystyle\sum_{i \in \mathcal{R}}
           w_i(f) \cdot r_i(m)$
        \Comment{Scaling factor: 4 muls + 3 adds}
    \State $\hat{t}_{\mathrm{cold}} \gets
           t_{\mathrm{ref}}(f) \cdot \sigma$
    \Statex
    \State $n \gets n(f, m)$
        \Comment{Observation count, 0 if unseen pair}
    \If{$n > 0$}
        \State $\hat{t} \gets \dfrac{n}{n + \kappa}
               \cdot \mathrm{EMA}(f,m)
               + \dfrac{\kappa}{n + \kappa}
               \cdot \hat{t}_{\mathrm{cold}}$
    \Else
        \State $\hat{t} \gets \hat{t}_{\mathrm{cold}}$
    \EndIf
    \Statex
    \State $\hat{t}_{\mathrm{pull}} \gets
           \Call{EstimatePullTime}{f, m}$
        \Comment{Section~\ref{sec:pulltime}}
    \State \Return $\hat{t} + \hat{t}_{\mathrm{pull}}$
\EndFunction

\Statex
\Statex

\Function{OnJobComplete}{$f, m, x$}
    \Comment{Called when observed runtime $x$ is reported}
    \State $n \gets n(f, m)$
    \If{$n = 0$}
        \State $\mathrm{EMA}(f,m) \gets \alpha \cdot x
               + (1-\alpha) \cdot \hat{t}_{\mathrm{cold}}(f,m)$
    \Else
        \State $\mathrm{EMA}(f,m) \gets \alpha \cdot x
               + (1-\alpha) \cdot \mathrm{EMA}(f,m)$
    \EndIf
    \State $n(f,m) \gets n + 1$
\EndFunction

\end{algorithmic}
\end{algorithm}
\bluecontentend

% =============================================================================
\section{Deficiencies of the Current Implementation}
\label{app:broken}
% =============================================================================
\bluecontentbegin

\begin{enumerate}[label=(\arabic*)]
    \item \textbf{Single-sample history.}
          \texttt{get\_latest\_run\_time} returns exactly one past runtime.
          A single anomalous run (GC pause, network hiccup) becomes the
          prediction for the next scheduling round. There is no smoothing.

    \item \textbf{MQTT round-trip latency.}
          The prediction fallback path requires a network round-trip to
          each provider. With MQTT pub/sub overhead and a 5\,s timeout,
          this adds \emph{seconds} of latency to every scheduling
          decision that involves an unseen $(f, m)$ pair.

    \item \textbf{Model retrained from scratch per request.}
          Every \texttt{PREDICT\_REQUEST} triggers a full
          \texttt{LinearRegression().fit()} on the provider. This is
          wasteful and introduces per-request latency.

    \item \textbf{Training data conflates all functions.}
          \texttt{eff\_score\_data.txt} mixes runs from different
          functions into one pool. A linear model fit on this data
          cannot distinguish a 55\,s CPU-heavy function from a 7\,s
          memory-heavy one --- the features (raw CPU/memory usage
          $\times$ efficiency) are not enough to separate them.

    \item \textbf{Efficiency scores are all 1.0.}
          In the current training data, every row has
          \texttt{cpu\_efficiency\_score = 1.0} and
          \texttt{memory\_efficiency\_score = 1.0}. The model has
          never seen heterogeneous machines, so the efficiency score
          features carry zero information.

    \item \textbf{Only two resource dimensions.}
          CPU and memory are tracked; disk I/O and network bandwidth
          are ignored. Functions with significant I/O dependence
          cannot be predicted accurately.

    \item \textbf{Pull time heuristic is baseless.}
          The factor $0.2$ applied to disk-cached pull time has no
          derivation and does not vary with image size or disk speed
          (Section~\ref{sec:pulltime}).

    \item \textbf{Hard-coded 1000\,ms default.}
          When all else fails, the scheduler uses 1000\,ms regardless
          of function or machine.
\end{enumerate}
\bluecontentend

% =============================================================================
\section{Computational Overhead Analysis}
\label{sec:overhead}
% =============================================================================
\bluecontentbegin

A primary design goal is that prediction overhead must be negligible
relative to the runtimes being predicted and the scheduling latency.

\subsection{Prediction Cost}

\begin{center}
\small
\begin{tabular}{@{}lll@{}}
    \toprule
    Operation & Cost & When \\
    \midrule
    Scaling factor $\sigma(f,m)$ & 4 muls + 3 adds & Every prediction \\
    Cold-start prediction & 1 mul & Every prediction \\
    Blending (if $n > 0$) & 2 muls + 1 add & If history exists \\
    EMA update & 2 muls + 1 add & On job complete \\
    \bottomrule
\end{tabular}
\end{center}

Total per prediction: $\sim$10 floating-point operations. On modern
hardware, this completes in \textbf{well under 1 microsecond}.

\subsection{Comparison with Current Implementation}

The current system dispatches an MQTT \texttt{PREDICT\_REQUEST} to
each provider when no historical data exists. This involves:
\begin{itemize}
    \item Network round-trip (MQTT pub/sub): $\sim$10--100\,ms.
    \item Provider-side LR training: $\sim$1--10\,ms.
    \item Timeout handling (up to 5\,s).
\end{itemize}

The proposed model eliminates this round-trip entirely. All data
(weights, reference runtimes, ratios) is preloaded at scheduler
startup. Predictions are computed locally in microseconds.

\subsection{Storage Cost}

\begin{center}
\small
\begin{tabular}{@{}llp{3.2cm}@{}}
    \toprule
    Data & Size & Stored Where \\
    \midrule
    Per-function benchmarks
      ($t_{\mathrm{ref}}, \mathbf{w}$)
      & $5 \!\times\! |\mathcal{F}|$ floats
      & Scheduler DB \\
    Per-machine ratios ($\mathbf{r}$)
      & $4 \!\times\! |\mathcal{M}|$ floats
      & User model \\
    Per-pair EMA state
      & $2 \!\times\! |\text{observed}|$
      & Summary table \\
    \bottomrule
\end{tabular}
\end{center}

For 100 functions and 50 machines: 500 + 200 + (at most 10,000)
floats. Entirely negligible.

\subsection{Offline Benchmarking Cost}

\begin{center}
\small
\begin{tabular}{@{}lll@{}}
    \toprule
    Task & Runs & Frequency \\
    \midrule
    Reference runtime & $B \!\times\! |\mathcal{F}|$ & Per function reg.\ \\
    Throttled runtimes & $4B' \!\times\! |\mathcal{F}|$ & Per function reg.\ \\
    Machine benchmarks & $4 \!\times\! |\mathcal{M}|$ & Per machine join \\
    \bottomrule
\end{tabular}
\end{center}

With $B = 5$, $B' = 3$: each new function requires $5 + 12 = 17$
runs on the BEM. Each new machine requires 4 benchmark runs. These
are one-time costs, amortized over all future invocations.
\bluecontentend

% =============================================================================
\section{Proposed Variant for Implementation}
\label{sec:variant}
% =============================================================================
\bluecontentbegin

Based on the structure of the current codebase and the analysis
above, we propose the following concrete implementation variant.

\subsection{Data Model Changes}

\begin{enumerate}
    \item \textbf{Extend the \texttt{User} (Provider) model} with two
          additional efficiency score fields:
          \texttt{disk\_efficiency\_score} and
          \texttt{network\_efficiency\_score}, stored as ratios
          relative to BEM (i.e., $r_{\texttt{disk}}$ and
          $r_{\texttt{net}}$). The existing
          \texttt{cpu\_efficiency\_score} and
          \texttt{memory\_efficiency\_score} already serve as
          $r_{\texttt{cpu}}$ and $r_{\texttt{mem}}$.

    \item \textbf{Extend the \texttt{Services} model} (or create a
          related \texttt{FunctionBenchmark} model) with fields:
          \texttt{ref\_runtime}, \texttt{w\_cpu}, \texttt{w\_mem},
          \texttt{w\_disk}, \texttt{w\_net}. Populated during function
          registration after BEM benchmarking.

    \item \textbf{Create a \texttt{PairRuntimeStats} model}:
          \begin{itemize}
              \item \texttt{provider} (FK $\to$ User)
              \item \texttt{service} (FK $\to$ Services)
              \item \texttt{ema\_runtime} (Float)
              \item \texttt{observation\_count} (Integer)
              \item Unique on (\texttt{provider}, \texttt{service})
          \end{itemize}
          Updated on every \texttt{JOB\_COMPLETE} event.
\end{enumerate}

\subsection{Prediction Flow}

Replace the current three-pass pipeline
(\texttt{db\_pass} $\to$ MQTT $\to$ \texttt{cache\_pass}) with:

\begin{enumerate}
    \item \textbf{Preload at scheduler startup:} Load all function
          benchmarks (weights, reference runtimes) and machine ratios
          into memory. Refresh on new registrations.

    \item \textbf{For each $(f, m)$ in cost matrix construction:}
          Compute $\hat{t}(f, m)$ via Algorithm~\ref{alg:predict}.
          All data is local --- no MQTT round-trip.

    \item \textbf{Apply pull time adjustment} via
          Equation~\ref{eq:pull-piecewise}:
          $V(f)/S_{\texttt{net}}$ or $V(f)/S_{\texttt{disk}}$
          from benchmarks, not the legacy
          $0.2 \times t_{\mathrm{pull}}^{\mathrm{last}}$ heuristic
          (Section~\ref{sec:pulltime}).

    \item \textbf{Feed the cost matrix to the ILP solver}
          (\texttt{minimize\_total\_cost}) as before.
\end{enumerate}

\subsection{EMA Update Path}

On receiving a \texttt{JOB\_COMPLETE} message from a provider:

\begin{enumerate}
    \item Extract \texttt{run\_time} from the completion payload.
    \item Upsert \texttt{PairRuntimeStats} for
          (\texttt{provider\_id}, \texttt{service\_id}):
          apply EMA update (Equation~\ref{eq:ema}).
    \item Next scheduling round picks up the updated EMA.
\end{enumerate}

\subsection{Eliminating the MQTT Prediction Path}

The MQTT \texttt{PREDICT\_REQUEST}/\texttt{PREDICT\_RESPONSE}
exchange exists solely because the scheduler lacks data to predict
locally for unseen pairs. With benchmark-derived weights and ratios
at the scheduler, this path is unnecessary. The provider-side
\texttt{handle\_predict\_request}, the scheduler-side
\texttt{\_fetch\_predicted\_runtimes\_mqtt\_batch}, and the
associated timeout/retry logic can all be retired.

\subsection{Summary of the Recommended Variant}

\begin{tcolorbox}[colback=blue!3, colframe=blue!50!black,
    title=Final Recommended Prediction Model]

\textbf{Cold-start} (no history for $(f, m)$):
\[
    \hat{t}(f, m) = t_{\mathrm{ref}}(f) \cdot
    \sum_{i \in \{\texttt{cpu}, \texttt{mem}, \texttt{disk},
    \texttt{net}\}}
    w_i(f) \cdot \frac{S_i(m_0)}{S_i(m)}
\]

\textbf{With history} ($n$ observations, EMA maintained online):
\[
    \hat{t}(f, m) = \frac{n}{n + \kappa} \cdot \mathrm{EMA}(f, m)
    + \frac{\kappa}{n + \kappa} \cdot \hat{t}_{\mathrm{cold}}(f, m)
\]

\textbf{Parameters:} $\alpha = 0.25$, $\kappa = 5$, $\theta = 0.5$,
$B = 5$, $B' = 3$.

\medskip

\textbf{Properties:}
\begin{itemize}[nosep]
    \item No ML model, no training, no inference overhead.
    \item Prediction in $< 1\,\mu$s (arithmetic only).
    \item Graceful transition from benchmarks to observations.
    \item Eliminates MQTT round-trip (seconds $\to$ $\mu$s).
    \item $O(1)$ storage and update per observed pair.
\end{itemize}
\end{tcolorbox}
\bluecontentend

% =============================================================================
\section{Limitations and Future Work}
% =============================================================================
\bluecontentbegin

\begin{enumerate}
    \item \textbf{Cross-resource coupling.} The additive scaling model
          assumes resources contribute independently to runtime.
          Functions with tightly interleaved CPU-memory access may
          exhibit non-linear interactions. Empirical validation will
          quantify this approximation error.

    \item \textbf{System noise under load.} Under heavy cluster load,
          contention for shared resources (cache, memory bus, network
          switch) may degrade performance beyond what static benchmarks
          predict. Incorporating real-time load signals as a
          multiplicative correction factor is a natural extension.

    \item \textbf{Benchmark staleness.} Machine performance may degrade
          over time (thermal throttling, disk aging). Periodic
          re-benchmarking or drift detection from EMA residuals
          ($|\mathrm{EMA} - \hat{t}_{\mathrm{cold}}|$ growing) can
          trigger re-calibration.

    \item \textbf{Network state variability.} Network bandwidth is the
          most volatile resource. Incorporating real-time network
          measurements could improve predictions for network-heavy
          workloads, though at considerable implementation complexity.
\end{enumerate}
\bluecontentend

\end{document}
